% © 2026 Ing. David Jaroš — CC BY-NC-ND 4.0
%
% This work is licensed under a Creative Commons Attribution-NonCommercial-NoDerivatives
% 4.0 International License (CC BY-NC-ND 4.0).
%
% License History: Earlier drafts (up to v0.3) were released under CC BY 4.0.
% From v0.4 onward, all material is released under CC BY-NC-ND 4.0 to protect
% the integrity of the theoretical work during ongoing academic development.
%
% See LICENSE.md for full license text.
%
% File: research_tracks/T3_ALPHA/bmod_from_action.tex
%
% Purpose: Track the derivation attempt of B_mod(p) = (p+1)/3 from S[Theta]
%   via winding-sector modular geometry.

% UBT-AI-PROVENANCE-BEGIN
% schema: ubt-ai-provenance/v1
% tier: C_working
% ai_assistance: disclosed
% human_review: risk-based
% editorial_responsibility: Ing. David Jaroš
% policy: ../../AI_PROVENANCE.md
% notice: Working material; exhaustive human review is not claimed.
% UBT-AI-PROVENANCE-END

\ifdefined\INCLUDEMODE
\else
  \documentclass[12pt,a4paper]{article}
  \usepackage[a4paper, margin=2.5cm]{geometry}
  \usepackage{amsmath, amssymb}
  \usepackage{hyperref}
  \usepackage{xcolor}
  \usepackage{tcolorbox}
  \tcbuselibrary{skins,breakable}

  \newtcolorbox{statusbox}[1][]{colback=gray!8!white,colframe=black!60,
    arc=3pt,boxrule=1pt,fonttitle=\bfseries,title=Status Summary,#1}

  \title{\textbf{Derivation Attempt of \(B_{\mathrm{mod}}(p)=\frac{p+1}{3}\) from \(S[\Theta]\)}}
  \author{Ing.\ David Jaro\v{s}}
  \date{2026-05-11}
  \begin{document}
  \maketitle
\fi

\section{Target statement}

The open Gap G137-B requires deriving
\begin{equation}
  B_{\mathrm{mod}}(p) = \frac{p+1}{3}
\end{equation}
from UBT dynamics, not as an external modular ansatz.

\section{Current structural bridge}

The arithmetic identity
\begin{equation}
  \frac{\mathrm{vol}(X_0(p))}{\pi}=\frac{\mu(\Gamma_0(p))}{3}=\frac{p+1}{3}
\end{equation}
is exact. The unresolved point is the physics map:
\[
S[\Theta]\ \Longrightarrow\ Z_{\mathrm{wind}}(p)\ \Longrightarrow\ B(p)
\]
with no fitted constants.

\section{Working UBT hypothesis}

For prime winding number \(n=p\), the winding Hilbert sector is expected to
organise over modular classes equivalent to \(\mathbb{P}^1(\mathbb{F}_p)\),
with cardinality \(p+1\). If the effective coefficient is the normalised class
count, then
\begin{equation}
  B = \frac{|\mathbb{P}^1(\mathbb{F}_p)|}{3} = \frac{p+1}{3}.
\end{equation}
The denominator 3 is interpreted as the existing UBT phase normalisation.

\section{Open technical obligations}

\begin{enumerate}
  \item Derive the \(\Gamma_0(p)\) (or equivalent) class action on the physical
    winding sector directly from \(S[\Theta]\).
  \item Show that the path-integral saddle sum yields a coefficient proportional
    to class count \(p+1\), not another weighting.
  \item Derive the exact \(1/3\) normalisation from the action measure and
    compactification conventions, with no adjustable prefactor.
\end{enumerate}

\begin{statusbox}
\textbf{Claim ``\(B_{\mathrm{mod}}(p)=\frac{p+1}{3}\) is derived from \(S[\Theta]\)''}:
\textbf{[OPEN]}.

\textbf{Arithmetic identity}: exact [L1 arithmetic layer].  
\textbf{UBT action-to-coefficient bridge}: unresolved [OPEN], currently usable only
as [MC] ansatz in self-consistency studies.
\end{statusbox}

\ifdefined\INCLUDEMODE
\else
  \end{document}
\fi
